
We now define a scheme-agnostic abstraction that captures the common structure
of the \GGM-tree mechanism across all schemes discussed in
Section~\ref{sec:background}.

\begin{definition}[\GZKST]
\sloppy
\label{def:gzkst}
A \emph{Generic ZK Seed Tree} (\GZKST) is a tuple
$\Pi = (\Setup,\, \Expand,\, \Chal,\, \Decide,\, \Publish,\, \Recon)$
parametrised by:
\begin{itemize}
  \item $\lam$: the security parameter;
  \item $t \in \mathbb{N}$: the number of parallel ZK repetitions;
  \item $\PRG : \{0,1\}^{\lam} \to \{0,1\}^{2\lam}$: a secure pseudorandom generator;
  \item $\calS$: the \emph{secret type}
        (e.g., permutations for LESS, matrix pairs for MEDS, party shares for MPCitH).
\end{itemize}
The six algorithms are:
\begin{enumerate}
  \item $\Setup(1^{\lam}) \to (\seed_r, \salt)$:
        samples a master seed $\seed_r \getsr \{0,1\}^{\lam}$ and a public salt.
  \item $\Expand(\seed_r, \salt) \to (\calT, (\eseed_0, \ldots, \eseed_{t-1}))$:
        builds a binary tree $\calT$ by top-down \PRG-expansion from $\seed_r$.
        Each leaf $i \in \{0, \ldots, t-1\}$ holds a seed $\eseed_i$.
        Tree topology may be complete or incomplete (scheme-dependent).
  \item $\Chal(\cmt, \msg) \to \mathbf{c} = (c_0, \ldots, c_{t-1})$:
        derives a Fiat--Shamir challenge from the commitment and message.
        $c_i = 0$ indicates that round $i$ requires the ephemeral seed;
        $c_i \neq 0$ indicates a secret-dependent response.
  \item $\Decide(\mathbf{c}) \to (\calH, \calR)$:
        partitions $\{0, \ldots, t-1\}$ into the hidden index set
        $\calH = \{i : c_i \neq 0\}$ and the revealed index set
        $\calR = \{i : c_i = 0\}$.
  \item $\Publish(\calT, \calH) \to \tpath$:
        constructs a flag structure $\calF : V(\calT) \to \{0,1\}$ and returns the
        minimal set of nodes $\tpath$ from which every $\eseed_i$ with
        $i \in \calR$ is reconstructible, while no $\eseed_i$ with
        $i \in \calH$ is derivable.
  \item $\Recon(\tpath, \mathbf{c}, \salt) \to \{\eseed_i : i \in \calR\}$:
        verifier-side reconstruction of revealed leaf seeds from $\tpath$.
\end{enumerate}
\end{definition}

% ------------------------------------------------------------
\subsection{The Flag Structure}
\label{sec:flag-structure}
% ------------------------------------------------------------

The flag structure $\calF$ produced internally by $\Publish$ is the key
intermediate object.  It is a labelling $\calF : V(\calT) \to \{0,1\}$ where:
\begin{itemize}
  \item $\calF(v) = 0$: node $v$ is \emph{publishable}---either
        directly included in $\tpath$ or derivable from a published ancestor;
  \item $\calF(v) = 1$: node $v$ is \emph{hidden}---it must not appear in
        $\tpath$ and must not be derivable from published nodes.
\end{itemize}

\noindent
The flag structure satisfies two invariants:

\begin{invariant}[Correctness]
\label{inv:correctness}
For every $i \in \calR$ there exists a node $v$ on the path from the root to
leaf $\eseed_i$ such that $\calF(v) = 0$ and $\calF(\mathsf{parent}(v)) = 1$.
\end{invariant}

\begin{invariant}[ZK hiding]
\label{inv:hiding}
For every $i \in \calH$, no ancestor of leaf $\eseed_i$ satisfies $\calF(v) = 0$.
\end{invariant}

The ZK property of the signature scheme is contingent on
Invariant~\ref{inv:hiding} holding for \emph{all} $i \in \calH$.  A fault that
changes $\calF(v)$ from $1$ to $0$ for any $v$ that is an ancestor of some
$\eseed_i$ with $i \in \calH$ \emph{breaks Invariant~\ref{inv:hiding}} and
may allow an attacker to reconstruct $\eseed_i$.

% ------------------------------------------------------------
\subsection{Mapping Schemes to the GZKST Model}
\label{sec:mapping}

We propose in this section to segment each of the ZK-based PQ signatures aforementionned into the elements of the GZKST tuple we defined in \ref{def:gzkst}.





% ------------------------------------------------------------



\subsubsection{LESS-V2}

\begin{itemize}
    \item $\Setup(1^{\lam}) \to (\seed_r, \salt)$ The root seed is sampled using the secret key seed. The salt is sampled uniformly at random.
    \item $\Expand(\seed_r, \salt) \to (\calT, (\eseed_0, \ldots, \eseed_{t-1}))$ 
LESS-v2 uses an \emph{incomplete} \GGM\ tree with $T = \lceil \log_2 t \rceil$
levels.  Leaves are distributed across the bottom two levels; the index offsets
are tracked via the arrays $\mathsf{npl}$, $\mathsf{ll}$, $\mathsf{off}$, and
$\mathsf{lsi}$.

    \item $\Chal(\cmt, \msg) \to \mathbf{c} = (c_0, \ldots, c_{t-1})$ In LESS, the commitment fed into the Fiat--Shamir hash is the concatenation
of all $t$ canonical-form matrices $\mathbf{B}^{(1)}, \ldots, \mathbf{B}^{(t)}$
together with the salt and the message.
The resulting digest \textsf{cmt} is then passed to \textsf{SampleChallenge}, which produces the challenge vector
$\mathbf{b} = (b^{(1)}, \ldots, b^{(t)}) \in \{0, \ldots, s-1\}^t$
of fixed weight $w$, where each non-zero entry $b^{(i)} \in \{1, \ldots, s-1\}$
indexes one of the $s-1$ public keys.
In the framework notation, the challenge value per round is
$c_i = b^{(i)}$: the condition $c_i = 0$ signals that round $i$ requires the
ephemeral seed $\eseed^{(i)}$, while $c_i \neq 0$ signals that a
secret-dependent coset representative must be provided instead.

    \item $\Decide(\mathbf{c}) \to (\calH, \calR)$ From the challenge vector $\mathbf{b}$, LESS computes the support
$U = \mathrm{supp}(\mathbf{b})$, i.e., the set of indices $i$ for
which $b^{(i)} \neq 0$. This directly induces the partition required by the framework.
    \item $\Publish(\calT, \calH) \to \tpath$  
    \paragraph{Flag tree.} LESS-v2 constructs an auxiliary
\emph{flag tree} of the same shape as the seed tree.  It is a labelling
$\calF : V(\calT) \to \{0,1\}$ where $\calF(v) = 0$ marks a publishable
node and $\calF(v) = 1$ marks a hidden node.  The flag tree is built in three
sequential phases:
\begin{enumerate}
  \item \textbf{Initialisation.}  Every node is set to $0$ except the root,
        which is set to $1$.
  \item \textbf{Leaf labelling} ($\LabelLeaves$).  For each leaf $i$, set
        \begin{equation*}
          \calF(\mathrm{leaf}_i) =
          \begin{cases} 1 & \text{if } b[i] \neq 0, \\ 0 & \text{otherwise,} \end{cases}
        \end{equation*}
        where $\mathbf{b} = (b[0], \ldots, b[t-1])$ is the Fiat--Shamir challenge
        vector of weight $w$ over the alphabet $\{0, \ldots, s-1\}$.
  \item \textbf{Upward propagation} ($\ComputeSeeds$).  For each internal node $v$
        with children $\ell$ and $r$:
        \begin{equation*}
          \calF(v) = \calF(\ell) \vee \calF(r).
        \end{equation*}
        This ensures the minimum set of nodes is published to reconstruct all
        revealed leaf seeds.
\end{enumerate}
The path is then constructed from the top down following this logic : if a node is flagged as true and its parent node is flagged as false, the seed in the current node is added to the path.

\item $\Recon(\tpath, \mathbf{c}, \salt) \to \{\eseed_i : i \in \calR\}$ 

On the verifier side, $\mathsf{PathToSeedTree}_t(h_0, \ldots, h_{t-1},\,
p,\, \alpha)$ reconstructs an empty tree of the same shape, places the
received path nodes at their correct positions, and re-expands each
downward via the PRG $\mathrm{G}$ until all $t$ leaf positions are
filled.
Positions $i \in \mathcal{H}$ are left empty (no seed on their
root-to-leaf path was ever published), while positions $i \in \mathcal{R}$
yield the recovered seeds $\sigma_i$, from which the verifier
re-derives $(\tilde{\mathbf{A}}_i, \tilde{\mathbf{B}}_i)$ and
re-computes $\tilde{\mathbf{G}}_i$ to verify the commitment $d$.

\end{itemize}




\begin{example}
Figure~\ref{fig:less-tree} illustrates an incomplete \GGM\ tree for $t = 6$ with
challenge $\mathbf{b} = (0,1,0,1,0,1)$, so
$\calH = \{1,3,5\}$ and $\calR = \{0,2,4\}$.
After all three phases, nodes $s_3$, $s_7$, and $s_9$ satisfy
$\calF(v)=0$ with $\calF(\mathsf{parent}(v))=1$ and thus enter $\tpath$,
giving the verifier access to $\eseed_0$, $\eseed_2$, and $\eseed_4$.
\end{example}

\begin{figure}[tb]
\centering
\begin{tikzpicture}[x=0.95cm, y=1.15cm]
  % ---- nodes (final state after all 3 phases) ----
  % Level 0
  \node[hidnode]  (s0)  at (3.5, 0)   {$s_0$};
  % Level 1
  \node[hidnode]  (s1)  at (1.5,-1)   {$s_1$};
  \node[hidnode]  (s2)  at (5.5,-1)   {$s_2$};
  % Level 2  (s3,s4 are leaves; s5,s6 are internal)
  \node[pathnode] (s3)  at (0.5,-2)   {$s_3$};
  \node[hidnode]  (s4)  at (2.5,-2)   {$s_4$};
  \node[hidnode]  (s5)  at (4.5,-2)   {$s_5$};
  \node[hidnode]  (s6)  at (6.75,-2)   {$s_6$};
  % Level 3  (leaves)
  \node[pathnode] (s7)  at (3.75,-3)  {$s_7$};
  \node[hidnode]  (s8)  at (5.25,-3)  {$s_8$};
  \node[pathnode] (s9)  at (6.00,-3)  {$s_9$};
  \node[hidnode]  (s10) at (7.25,-3)  {$s_{10}$};
  % ---- ephemeral-seed labels below leaves ----
  \node[leaflbl] at (0.5,-2.55)  {$\eseed_0$};
  \node[leaflbl] at (2.5,-2.55)  {$\eseed_1$};
  \node[leaflbl] at (3.75,-3.55) {$\eseed_2$};
  \node[leaflbl] at (5.25,-3.55) {$\eseed_3$};
  \node[leaflbl] at (5.75,-3.55) {$\eseed_4$};
  \node[leaflbl] at (7.25,-3.55) {$\eseed_5$};
  % ---- edges ----
  \draw[treeedge] (s0)--(s1); \draw[treeedge] (s0)--(s2);
  \draw[treeedge] (s1)--(s3); \draw[treeedge] (s1)--(s4);
  \draw[treeedge] (s2)--(s5); \draw[treeedge] (s2)--(s6);
  \draw[treeedge] (s5)--(s7); \draw[treeedge] (s5)--(s8);
  \draw[treeedge] (s6)--(s9); \draw[treeedge] (s6)--(s10);
  % ---- "incomplete" annotation ----
  \node[font=\tiny,gray] at (3.5,-2.55) {leaf level~2};
  % ---- legend ----
  \matrix[draw=gray!40, rounded corners=2pt, inner sep=4pt,
          above left=0.2cm and -0.8cm of s1,
          matrix of nodes, nodes={font=\tiny, inner sep=2pt},
          column sep=4pt, row sep=1pt] {
    \node[hidnode,  minimum size=8pt]{}; & \node{$\calF=1$ (hidden)};  \\
    \node[pathnode, minimum size=8pt]{}; & \node{$\calF=0$, in $\tpath$}; \\
  };
\end{tikzpicture}
\caption{LESS-v2 incomplete \GGM\ flag tree, $t=6$, after all three construction
         phases.  Challenge $\mathbf{b}=(0,1,0,1,0,1)$ hides rounds
         $\calH=\{1,3,5\}$.  Path nodes (blue, thick border) give the verifier
         $\eseed_0$, $\eseed_2$, $\eseed_4$.  Injecting a fault on any red node
         exposes the corresponding hidden seed alongside its response, yielding a
         dual revelation.}
\label{fig:less-tree}
\end{figure}


Table~\ref{tab:less-params} lists the parameter sets of LESS.

\begin{table}[tb]
\centering
\caption{LESS-v2 parameter sets~\cite{LESS-spec}.}
\label{tab:less-params}
\begin{tabular}{@{}lcccccc@{}}
\toprule
Parameter set   & $n$ & $k$ & $q$ & $t$ & $w$ & $s$ \\
\midrule
LESS-252-45     & 252 & 126 & 127 &  45 & 34 & 8 \\
LESS-252-68     & 252 & 126 & 127 &  68 & 42 & 4 \\
LESS-252-192    & 252 & 126 & 127 & 192 & 36 & 2 \\
LESS-400-102    & 400 & 200 & 127 & 102 & 61 & 4 \\
LESS-400-220    & 400 & 200 & 127 & 220 & 68 & 2 \\
LESS-548-137    & 548 & 274 & 127 & 137 & 79 & 4 \\
LESS-548-345    & 548 & 274 & 127 & 345 & 75 & 2 \\
\bottomrule
\end{tabular}
\end{table}

\subsubsection{MEDS}

\begin{itemize}
    \item $\Setup(1^{\lam}) \to (\seed_r, \salt)$ 
    
    MEDS samples a fresh $\ell_{\mathrm{sec\_seed}}$-byte random scalar, then derives both the master tree seed
$\rho \in \mathcal{B}^{\ell_{\mathrm{tree\_seed}}}$ and the public salt
$\alpha \in \mathcal{B}^{\ell_{\mathrm{salt}}}$ in a single XOF call. 
    \item $\Expand(\seed_r, \salt) \to (\calT, (\eseed_0, \ldots, \eseed_{t-1}))$ 
    
    From $(\rho, \alpha)$, MEDS calls $\mathsf{SeedTree}_t(\rho, \alpha)$
(line~13) to build a complete binary tree of height
$\lceil \log_2 t \rceil$ and return its first $t$ leaves
$\sigma_0, \ldots, \sigma_{t-1} \in \mathcal{B}^{\ell_{\mathrm{tree\_seed}}}$.
    \item $\Chal(\cmt, \msg) \to \mathbf{c} = (c_0, \ldots, c_{t-1})$ 
    
    Each leaf $\sigma_i$ is then expanded into three sub-seeds
$(\sigma_{\tilde{A}_i},\, \sigma_{\tilde{B}_i},\, \sigma_i')$ with $\mathrm{XOF}$.
from which ephemeral invertible matrices sampled then generator matrixes computed.
The commitment is formed by hashing the compressed systematic forms of all
$t$ ephemeral generator matrices $\tilde{\mathbf{G}}_0, \ldots,
\tilde{\mathbf{G}}_{t-1}$ together with the message $m$.

The digest $d$ is then parsed by $\mathsf{ParseHash}_{s,t,w}(d)$
to produce the challenge vector
$h_0, \ldots, h_{t-1} \in \{0, \ldots, s-1\}$
of fixed weight $w$.

In the framework notation, $c_i = h_i$. 

    \item $\Decide(\mathbf{c}) \to (\calH, \calR)$ The condition $c_i = 0$ signals
that round $i$ requires the ephemeral seed $\mathsf{es}_i$, while
$c_i \neq 0$ indicates that a secret-dependent matrix response must be
provided.
    \item $\Publish(\calT, \calH) \to \tpath$  
    
    MEDS calls $\mathsf{SeedTreeToPath}_t(h_0, \ldots, h_{t-1},\, \rho,\,
\alpha)$, which rebuilds the full seed tree, then for each
$i \in \mathcal{H}$ removes the label of leaf $i$ and propagates the
deletion upward, erasing every ancestor whose subtree is entirely hidden.
The function then serialises in depth-first order all remaining non-empty
node labels at the highest possible level, zero-padded to
$\ell_{\mathrm{path}}$ bytes, yielding the seed path $p$.

\item $\Recon(\tpath, \mathbf{c}, \salt) \to \{\eseed_i : i \in \calR\}$ 

On the verifier side, $\mathsf{PathToSeedTree}_t(h_0, \ldots, h_{t-1},\,
p,\, \alpha)$ reconstructs an empty tree of the same shape, places the
received path nodes at their correct positions, and re-expands each
downward via the PRG $\mathrm{G}$ until all $t$ leaf positions are
filled.
Positions $i \in \mathcal{H}$ are left empty (no seed on their
root-to-leaf path was ever published), while positions $i \in \mathcal{R}$
yield the recovered seeds $\sigma_i$, from which the verifier
re-derives $(\tilde{\mathbf{A}}_i, \tilde{\mathbf{B}}_i)$ and
re-computes $\tilde{\mathbf{G}}_i$ to verify the commitment $d$.
\end{itemize}


\begin{example}
Figure~\ref{fig:meds-tree} shows a complete \GGM\ tree for $t = 4$ with
challenge $\mathbf{h} = (0,1,0,1)$, so $\calH = \{1,3\}$ and $\calR = \{0,2\}$.
Nodes $s_3$ and $s_5$ enter $\tpath$.
\end{example}

\begin{figure}[tb]
\centering
\begin{tikzpicture}[x=1.1cm, y=1.15cm]
  % Level 0
  \node[hidnode]  (s0) at (2,  0)  {$s_0$};
  % Level 1
  \node[hidnode]  (s1) at (0.5,-1) {$s_1$};
  \node[hidnode]  (s2) at (3.5,-1) {$s_2$};
  % Level 2 (leaves)
  \node[pathnode] (s3) at (-0.5,-2){$s_3$};
  \node[hidnode]  (s4) at ( 1.5,-2){$s_4$};
  \node[pathnode] (s5) at ( 2.5,-2){$s_5$};
  \node[hidnode]  (s6) at ( 4.5,-2){$s_6$};
  % Ephemeral seed labels
  \node[leaflbl] at (-0.5,-2.55){$\eseed_0$};
  \node[leaflbl] at ( 1.5,-2.55){$\eseed_1$};
  \node[leaflbl] at ( 2.5,-2.55){$\eseed_2$};
  \node[leaflbl] at ( 4.5,-2.55){$\eseed_3$};
  % Edges
  \draw[treeedge] (s0)--(s1); \draw[treeedge] (s0)--(s2);
  \draw[treeedge] (s1)--(s3); \draw[treeedge] (s1)--(s4);
  \draw[treeedge] (s2)--(s5); \draw[treeedge] (s2)--(s6);
  % Legend
  \matrix[draw=gray!40, rounded corners=2pt, inner sep=4pt,
          above right=0.6cm and -1cm of s1,
          matrix of nodes, nodes={font=\tiny, inner sep=2pt},
          column sep=4pt, row sep=1pt] {
    \node[hidnode,  minimum size=8pt]{}; & \node{$\calF=1$};  \\
    \node[pathnode, minimum size=8pt]{}; & \node{in $\tpath$}; \\
  };
\end{tikzpicture}
\caption{MEDS complete \GGM\ flag tree, $t=4$.
         Challenge $\mathbf{h}=(0,1,0,1)$ hides $\calH=\{1,3\}$.
         Path nodes $s_3$ and $s_5$ (blue) give the verifier $\eseed_0$ and
         $\eseed_2$.  A fault on $s_4$ or $s_6$ exposes the corresponding hidden
         seed alongside its ZK response.}
\label{fig:meds-tree}
\end{figure}

\begin{table}[tb]
\centering
\caption{MEDS parameter sets~\cite{ChouNPRRST23}.}
\label{tab:meds-params}
\begin{tabular}{@{}lccccc@{}}
\toprule
Parameter set & $n$ & $k$ & $q$ & $t$ & $s$ \\
\midrule
MEDS-9923b    & 14 &  7 & 509 & 192 & 2 \\
MEDS-13220b   & 22 & 11 & 127 & 192 & 2 \\
MEDS-41711b   & 30 & 15 & 509 & 392 & 2 \\
\bottomrule
\end{tabular}
\end{table}

\subsubsection{CROSS}

\begin{itemize}
    \item $\Setup(1^{\lam}) \to (\seed_r, \salt)$: The root seed is of size $\lam$ and the salt $2\lam$. They are both theoretically drawn uniformly at random.
    \item $\Expand(\seed_r, \salt) \to (\calT, (\eseed_0, \ldots, \eseed_{t-1}))$: 

CROSS-v1 uses a complete \GGM\ tree identical in structure to LESS-v1.
CROSS-v2 may adopt an optimised variant; the reference implementation uses a
$\SeedTreePaths$ mechanism structurally equivalent to that of MEDS.

The GGM tree in CROSS is a full binary tree where each parent node's $\lam$-bit seed is concatenated with a global salt and the node's unique index, then processed through a CSPRNG (SHAKE) to deterministically produce two $\lam$-bit strings for its left and right children.
      \item $\Chal(\cmt, \msg) \to \mathbf{c} = (c_0, \ldots, c_{t-1})$ 
      
      CROSS applies the Fiat--Shamir transform in two successive rounds.

First, the $t$ commitment pairs $(\mathsf{cmt}_0[i], \mathsf{cmt}_1[i])$,
the salt, and the message are hashed to produce the first challenge vector
$\mathsf{chall}_1 \in (\mathbb{F}_p^*)^t$, from which each first response
$\mathsf{digest}_{y_i} = \mathrm{Hash}({\bf u}'_i + \mathsf{chall}_1[i]
\cdot {\bf e}'_i)$ is derived.

A second hash over all first-challenge digests then yields the
second challenge vector $\mathsf{chall}_2 \in \mathcal{B}(t, w)$ of fixed Hamming weight $w$.
In the framework notation, $c_i = \mathsf{chall}_2[i] \in \{0,1\}$: the
condition $c_i = 0$ signals that round $i$ requires the ephemeral seed
$\mathsf{es}_i$, while $c_i = 1$ signals that a secret-dependent response
$({\bf y}_i, {\bf v}_i)$ must be provided.

  \item $\Decide(\mathbf{c}) \to (\calH, \calR)$

    The challenge vector determines which of the $t$ ephemeral seeds to hide:
$\calH = \{i : c_i \neq 0\}$ and $\calR = \{i : c_i = 0\}$. 

  \item $\Publish(\calT, \calH) \to \tpath$  
  
CROSS computes the seed path $\mathsf{Path}$ by applying the seed-tree
path algorithm to the full tree with respect to $\mathcal{H}$: for each
$i \in \mathcal{H}$, the label of leaf $i$ is removed and all ancestor
labels on its root-to-leaf path are deleted, then the minimal set of
surviving boundary nodes is serialised in depth-first order.
This corresponds to the framework flag structure
$\mathcal{F} : V(\mathcal{T}) \to \{0,1\}$ with $\mathcal{F}(v) = 1$
marking deleted nodes and $\tpath$ being the remaining boundary.

Additionally, CROSS computes a \emph{Merkle proof} $\mathsf{Proof}$ over
the $t$ commitments $\mathsf{cmt}_0[0], \ldots, \mathsf{cmt}_0[t-1]$,
authenticating the $w$ values $\mathsf{cmt}_0[i]$ for $i \in \mathcal{H}$
that the verifier cannot recompute from any revealed seed.

This Merkle proof is a CROSS-specific mechanism with no analogue in LESS
or MEDS, arising from the dual-commitment structure of the underlying
5-pass protocol.
  \item $\Recon(\tpath, \mathbf{c}, \salt) \to \{\eseed_i : i \in \calR\}$ 
  
  On the verifier side, the received $\mathsf{Path}$ nodes are placed onto
an empty tree of the same structure and re-expanded downward via the CSPRNG
until all $t$ leaf positions are populated.
Positions $i \in \mathcal{H}$ remain empty, as no seed on their
root-to-leaf path was ever published; the verifier instead uses the
explicit response $\mathsf{resp}[i] = ({\bf y}_i, {\bf v}_i)$ together
with the Merkle proof to authenticate $\mathsf{cmt}_0[i]$ and verify the
round.

Positions $i \in \mathcal{R}$ yield the recovered seeds $\mathsf{Seed}[i]$,
from which the verifier re-derives $({\bf e}'_i, {\bf u}'_i)$,
recomputes $\mathsf{cmt}_1[i]$, and checks consistency with the signed
commitment digest.

\end{itemize}
\begin{example}
Figure~\ref{fig:cross-tree} shows a complete \GGM\ tree for $t = 4$ with ternary
challenge $\mathbf{c} = (0,\, 2,\, 0,\, 1)$.  Because $c_1 = 2 \neq 0$ and
$c_3 = 1 \neq 0$, the hidden set is $\calH = \{1,3\}$ and the path nodes are
$s_3$ and $s_5$.  Faulting $s_4$ exposes $\eseed_1$ and thereby leaks the
\emph{second} secret component (corresponding to $c_1 = 2$);
faulting $s_6$ leaks the \emph{first} (corresponding to $c_3 = 1$).
\end{example}

\begin{figure}[tb]
\centering
\begin{tikzpicture}[x=1.1cm, y=1.15cm]
  % Level 0
  \node[hidnode]  (s0) at (2,   0) {$s_0$};
  % Level 1
  \node[hidnode]  (s1) at (0.5,-1) {$s_1$};
  \node[hidnode]  (s2) at (3.5,-1) {$s_2$};
  % Level 2 (leaves)
  \node[pathnode] (s3) at (-0.5,-2){$s_3$};
  \node[hidnode]  (s4) at ( 1.5,-2){$s_4$};
  \node[pathnode] (s5) at ( 2.5,-2){$s_5$};
  \node[hidnode]  (s6) at ( 4.5,-2){$s_6$};
  % Ephemeral-seed labels + challenge values
  \node[leaflbl] at (-0.5,-2.6){$\eseed_0,\; c_0{=}0$};
  \node[leaflbl] at ( 1.5,-2.6){$\eseed_1,\; c_1{=}2$};
  \node[leaflbl] at ( 2.5,-2.6){$\eseed_2,\; c_2{=}0$};
  \node[leaflbl] at ( 4.5,-2.6){$\eseed_3,\; c_3{=}1$};
  % Edges
  \draw[treeedge] (s0)--(s1); \draw[treeedge] (s0)--(s2);
  \draw[treeedge] (s1)--(s3); \draw[treeedge] (s1)--(s4);
  \draw[treeedge] (s2)--(s5); \draw[treeedge] (s2)--(s6);
  % Annotations on hidden leaves
  \node[font=\tiny,red!70!black] at (1.5,-2) {\rotatebox{0}{}};
  % Legend
  \matrix[draw=gray!40, rounded corners=2pt, inner sep=4pt,
          above right=0.2cm and -2.5cm of s1,
          matrix of nodes, nodes={font=\tiny, inner sep=2pt},
          column sep=4pt, row sep=1pt] {
    \node[hidnode,  minimum size=8pt]{}; & \node{$\calF=1$, leaks $S_{c_i}$};  \\
    \node[pathnode, minimum size=8pt]{}; & \node{$\calF=0$, in $\tpath$}; \\
  };
\end{tikzpicture}
\caption{CROSS complete \GGM\ flag tree, $t=4$, ternary challenge
         $\mathbf{c}=(0,2,0,1)$.  Faulting hidden node $s_4$ exposes $\eseed_1$
         and leaks the second secret component ($c_1=2$); faulting $s_6$ leaks
         the first ($c_3=1$).  The ternary alphabet means that different hidden
         leaves point to \emph{different} components of the long-term secret.}
\label{fig:cross-tree}
\end{figure}

\subsubsection{MPCitH / VOLEitH Family }

\begin{itemize}
    \item $\Setup(1^{\lam}) \to (\seed_r, \salt)$ TODO
    \item $\Expand(\seed_r, \salt) \to (\calT, (\eseed_0, \ldots, \eseed_{t-1}))$ TODO
    \item $\Chal(\cmt, \msg) \to \mathbf{c} = (c_0, \ldots, c_{t-1})$ TODO
    \item $\Decide(\mathbf{c}) \to (\calH, \calR)$ TODO
    \item $\Publish(\calT, \calH) \to \tpath$  TODO
\end{itemize}
